% ---------------------------------------------------------------------------------------------
%  Overview of the argument (hand-written). Input by handout.tex, after its heading of the same name.
%  A fragment: no preamble. Every boxed statement was compared with its Lean statement at the pinned
%  commit; Lean names and file names stand in the closing table only (rows made from reading/census.json).
% ---------------------------------------------------------------------------------------------
\providecommand{\ovsec}[2]{\needspace{5\baselineskip}\vspace{0.8em}\phantomsection
  \addcontentsline{toc}{subsection}{#1. #2}{\sffamily\bfseries\large\color{ink} #1.\quad #2\par}\vspace{0.1em}}
\providecommand{\ovthm}[2]{\begin{tcolorbox}[colback=boxbg,colframe=accent,boxrule=0.6pt,arc=2pt,left=3.5mm,
  right=3.5mm,top=1.8mm,bottom=1.8mm,before skip=0.8em,after skip=0.8em]\setlength{\parskip}{0.35em}%
  {\sffamily\bfseries\color{accent}#1.}\enspace #2\end{tcolorbox}}
\providecommand{\ovstm}[2]{\begin{tcolorbox}[colback=boxbg,colframe=boxbg,boxrule=0pt,arc=2pt,left=3.5mm,
  right=3.5mm,top=1.8mm,bottom=1.8mm,before skip=0.8em,after skip=0.8em]\setlength{\parskip}{0.35em}%
  {\sffamily\bfseries #1.}\enspace #2\end{tcolorbox}}
\providecommand{\ovset}[2]{\begin{tcolorbox}[breakable,colback=boxbg,colframe=boxbg,boxrule=0pt,arc=2pt,left=3.5mm,
  right=3.5mm,top=1.8mm,bottom=1.8mm,before skip=0.8em,after skip=0.8em]\small\setlength{\parskip}{0.3em}%
  {\sffamily\bfseries #1.}\enspace #2\end{tcolorbox}}
\providecommand{\ovhead}[1]{\textbf{#1}\enspace}
\providecommand{\ovRe}{\operatorname{Re}}
\providecommand{\ovIm}{\operatorname{Im}}
\providecommand{\ovord}{\operatorname{ord}}
\providecommand{\ovres}{\operatorname{res}}
\providecommand{\ovAut}{\operatorname{Aut}(\mathbb{C})}
\providecommand{\ovRow}{\mathrm{Row}}
\providecommand{\ovCol}{\mathrm{Col}}
\providecommand{\ovColram}{\mathrm{Col}^{\mathrm{ram}}}
\providecommand{\ovrem}{\mathrm{rem}}
\providecommand{\ovhmain}{h_{\mathrm{main}}}
\providecommand{\ovbudget}{\mathrm{budget}}
\providecommand{\ovtops}{\mathrm{tops}}
\providecommand{\ovthick}{\mathrm{thick}}
\providecommand{\ovAsk}{\mathrm{Ask}}
\providecommand{\ovG}{\mathsf{G}}
\providecommand{\ovU}{\mathsf{U}}
\providecommand{\ovP}{\mathcal{P}}
\providecommand{\ovA}{\mathcal{A}}
\providecommand{\ovT}{\mathcal{T}}
\providecommand{\ovcov}{\operatorname{covol}}
\providecommand{\ovto}[1]{\ \overset{\text{\S#1}}{\Longrightarrow}\ }

\ovsec{1}{What is claimed}

The formal text ends with five theorems. Let $\zeta$ be the Riemann zeta function, $M\ge 1$ an integer, $\chi$ a
Dirichlet character modulo $M$, primitive or not, and $L(s,\chi)$ its $L$-function. The five theorems say:
\begin{enumerate}[label=(\arabic*),leftmargin=2.3em,itemsep=0.05em,topsep=0.15em]
\item every $s\in\mathbb{C}$ with $\zeta(s)=0$, $s\notin\{-2,-4,-6,\dots\}$ and $s\ne 1$ has $\ovRe s=\tfrac12$; this is
  the proposition that Mathlib defines as the Riemann hypothesis;
\item the same statement, written out in the statement file of the repository; its formal proof is theorem (1)
  itself, so the two propositions agree by definition;
\item $L(s,\chi)\ne 0$ whenever $\ovRe s>\tfrac12$, unless $\chi$ is principal and $s=1$;
\item every zero $s$ of $L(\,\cdot\,,\chi)$ with $0<\ovRe s<1$ has $\ovRe s=\tfrac12$;
\item statement (4) for primitive $\chi$. It is a special case of (4): its proof does not use primitivity.
\end{enumerate}
In Mathlib $\zeta$ and $L(\,\cdot\,,\chi)$ are functions on all of $\mathbb{C}$: the meromorphic continuations, with a
conventional finite value at the pole $s=1$ of $\zeta$ and of $L(\,\cdot\,,\chi)$ for principal $\chi$. This is why
$s=1$ is excluded by hand. A Dirichlet character modulo $M$ is a multiplicative map
$\mathbb{Z}/M\to\mathbb{C}$ that vanishes off the units, and the principal character is $1$ on the units. An imprimitive
character lacks finitely many Euler factors of the primitive one, which adds zeros on the line $\ovRe s=0$ only; hence
the open strip in (4), and (3)--(5) say nothing about $\ovRe s\le 0$.

\ovhead{What one must trust.} Three things, and no more: the statement file, which imports only Mathlib and defines
the propositions (2)--(5); Mathlib's definitions of $\zeta$, of Dirichlet characters and of $L(s,\chi)$, at the pinned
version; and Lean's kernel with its three standard axioms (propositional extensionality, choice, soundness of
quotients). Every other definition of the repository, and so every object of this overview, occurs only inside proofs.

\ovhead{Status.} The authors report that all files compile with no proof omitted and that the kernel accepts the
five theorems from these axioms; the record of checks at the end of this handout says who checked what, including
how far our own replay has got. That is a fact about a formal text. No mathematician has read the argument through, and no paper
accompanies it. The result is therefore a \emph{claim}: this handout describes what the formal proof asserts and
how it argues, and it does not say that the Riemann hypothesis is proved. The repository attributes the proof to Levent
Alpöge, with Claude, and says that its Lean text was written by Claude; the compact edition read here rewrites an
earlier, longer formal proof and follows its method.

\ovhead{Three grades of statement.} \emph{Kernel-checked, as reported}: every statement of the formal text, so every
theorem in a box below. \emph{Read}: rewritten by us as mathematics from structured proofs that a program compared with the source;
the chapter list after this overview shows the state of each unit. \emph{Quoted}: results that the formal
text proves in its background library and that we state without reading their proofs: classical results, cited
from the literature, and one result that the literature covers only in part (Section~6). In this overview every
theorem in a box was compared, hypothesis by hypothesis, with its Lean statement, and the closing table gives the
places; the objects of Setting~2 are described after their Lean definitions and the authors' comments, and the
chapters define them. The deductions carried out in Sections 2 to 4 follow short Lean proofs that were read in the
source. The paragraphs headed
\emph{Idea} are not readings of proofs: they report the authors' comments and the statements of intermediate lemmas,
and the chapters will replace them. The argument runs as follows; A, B, C are the three main theorems:
\[
\text{A, B, C}\ovto{4}\text{main inequality}\ovto{3}\text{one-sided bound for primes}\ovto{2}\text{(3)}\ovto{2}\text{(1)--(5)}.
\]

\ovsec{2}{From primes to zeros}

\ovhead{The weighted prime sum.} Let $d$ be a prime and $M\ge1$ an integer. Put $K=\mathbb{Q}(\sqrt{-d})$, and let
$\chi_K$ be the quadratic character of the field $\mathbb{Z}/d$, a Dirichlet character modulo $d$ (the Legendre
symbol $(\,\cdot\,/d)$ when $d$ is odd); for $d\equiv 3 \bmod 4$ it is the quadratic character of $K$. With
$\Lambda$ the von Mangoldt function, define for integers $\nu\ge1$
\[
\lambda_{d,M}(\nu)=\Lambda(\nu)\bigl(1-\chi_K(\nu)\bigr)\ \ \text{if}\ \nu\equiv1 \bmod M,\qquad \lambda_{d,M}(\nu)=0\ \ \text{otherwise.}
\]
So $\lambda_{d,M}\ge0$, and, when $d\equiv3\bmod4$ (as for every $d$ used below, all $\equiv7\bmod8$), $\lambda_{d,M}$
vanishes at the powers of the primes that split in $K$. For real $x$ the
\emph{remainder} is
\begin{gather*}
R_{d,M}(x)=\sum_{1\le\nu\le x}\lambda_{d,M}(\nu)\Bigl(\frac1\nu-\frac1x\Bigr)-\frac{\log x}{\varphi(M)}-\kappa_{d,M},\\
\kappa_{d,M}=\frac{1}{\varphi(M)}\,\ovRe\Bigl(\sum_{\xi}\Bigl(\frac{\tilde L'}{\tilde L}(1,\xi\chi_K)-\frac{\tilde L'}{\tilde L}(1,\xi)\Bigr)-1\Bigr).
\end{gather*}
Here $\varphi$ is Euler's function, $\xi$ runs over the Dirichlet characters modulo $M$, $\xi\chi_K$ is the
character modulo $Md$ with the values $\xi(\nu)\chi_K(\nu)$, and $\tilde L(s,\xi)$ is the
entire function $L(s,\xi)$, or $(s-1)L(s,\xi)$ if $\xi$ is principal. The term $\log x/\varphi(M)+\kappa_{d,M}$ is the
expected value of the sum. (A remark of ours: by the explicit formula, the generalized Riemann hypothesis gives
$R_{d,M}(x)=O(x^{-1/2})$.)

\ovhead{Zeros as poles.} Let the prime $d$ be $\equiv7\bmod8$ and prime to $M$. The Laplace transform
$\Phi_{d,M}(s)=\int_0^\infty R_{d,M}(e^{v})\,e^{(1-s)v}\,dv$, defined for $\ovRe s>2$, continues to a meromorphic function
on $\mathbb{C}$ that is analytic at every point of the closed half-plane $\ovRe s\ge1$, the point $s=1$ included:
the constant $\kappa_{d,M}$ is the one that cancels the pole at $s=1$. In $\ovRe s>0$ its poles are at most simple,
and for $\ovRe\rho>0$, $\rho\ne1$, the product $(s-\rho)\Phi_{d,M}(s)$ tends, as $s\to\rho$, to
\[
\ovres_{d,M}(\rho)=\frac{1}{\varphi(M)\,\rho(1-\rho)}\sum_{\xi}\Bigl(\ovord_{\rho}\tilde L(\,\cdot\,,\xi)-\ovord_{\rho}\tilde L(\,\cdot\,,\xi\chi_K)\Bigr),
\]
where $\ovord_{\rho}f$ is the order of vanishing at $\rho$ of an entire function $f\ne0$; one puts
$\ovres_{d,M}(1)=0$, the limit at $\rho=1$. A zero of some $L(\,\cdot\,,\xi)$, $\xi$ modulo $M$, counts positively, and a zero of a twist $L(\,\cdot\,,\xi\chi_K)$
negatively.

\ovstm{One-sided criterion}{Let $M\ge1$ be an integer, $c$ a unit of $\mathbb{Z}/M$ and $\beta>\tfrac12$ a real number.
Suppose that every prime $d\equiv c\bmod M$ satisfies $d\equiv 7\bmod 8$, is prime to $M$, and has
\[R_{d,M}(x)\le x^{\beta-1}\quad\text{for all sufficiently large real }x.\]
Then $\tilde L(\rho,\chi)\ne0$ for every Dirichlet character $\chi$ modulo $M$ and every $\rho\in\mathbb{C}$ with
$\ovRe\rho>\beta$.}

The proof has two inputs from the background library; one may assume $\ovRe\rho<1$, the rest being classical.
Given $\rho$, the first input provides one prime $d\equiv c\bmod M$ for which no
$\tilde L(s,\xi\chi_K)$, $\xi$ modulo $M$, vanishes in the box $\ovRe s\ge\beta$, $|\ovIm s|\le|\ovIm\rho|$; it rests
on estimates of zero-density type for quadratic twists, and Section~6 says what of it the literature covers. This is why
the bound is asked of every prime of the class: the prime depends on the height of $\rho$. For this $d$ the numbers
$\ovres_{d,M}(s)$ at real points $s>\beta$ are non-negative, and a constant multiple of $e^{(\beta-1)v}$ exceeds
$R_{d,M}(e^{v})$ for all $v\ge 0$. Now Landau's theorem applies: the Laplace transform of a non-negative function has a singularity
at the real point of its abscissa of convergence. This is the integral form of his theorem on Dirichlet series with
non-negative coefficients, used in a version that tolerates real poles of the favourable sign. So $\Phi_{d,M}$ has
no pole in $\ovRe s>\beta$, and $\ovres_{d,M}(\rho)=0$. At $\rho$ the twisted terms vanish by the choice of $d$; the
remaining terms are non-negative with sum zero, so each is zero.

\ovstm{One-sided bound}{Let $N_0$ be a positive multiple of $18480$, let $d$ be a prime with
$d\equiv N_0/2-1\bmod N_0$, and let $\varepsilon$ be real with $0<\varepsilon<\tfrac1{100}$. Then
\[R_{d,N_0}(x)\le x^{8\varepsilon-1/2}\quad\text{for all sufficiently large real }x.\]}

\ovhead{How the five theorems follow.} Let $\chi$ be a character modulo $M$ and suppose $L(\rho,\chi)=0$ with
$\ovRe\rho>\tfrac12$, the pole excepted. Put $N_0=18480\,M$ and $c=N_0/2-1$. Then $c^2\equiv1\bmod N_0$, so $c$ is a
unit, and every prime $d\equiv c\bmod N_0$ is prime to $N_0$ and is $\equiv7\bmod 8$, because $16$ divides
$18480=2^4\cdot3\cdot5\cdot7\cdot11$. The character modulo $N_0$ induced by $\chi$ still has
$\tilde L(\rho,\,\cdot\,)=0$, the two $L$-functions differing by Euler factors without zeros in $\ovRe s>0$. Choose
$\varepsilon=(\min(\ovRe\rho,1)-\tfrac12)/51$, so that $\varepsilon<\tfrac1{100}$ and
$\beta=8\varepsilon+\tfrac12<\ovRe\rho$. The one-sided bound, for every such $d$, is the hypothesis of the criterion
with modulus $N_0$, and the criterion contradicts the zero. This is theorem (3). The functional equation, taken from
Mathlib for primitive characters and extended to all characters through their missing Euler factors, pairs a zero $s$
in the strip with the zero $1-s$ of $L(\,\cdot\,,\bar\chi)$, which gives (4) and (5); for $M=1$, with the trivial zeros
and the pole treated by the functional equation of $\zeta$, it gives (1) and (2). Everything else in the formal text
serves the one-sided bound.

\ovsec{3}{The main inequality}

\ovset{Setting 1 (parameters and data)}{\emph{Parameters} are a pair $P=(N_0,d)$: a positive multiple $N_0$ of $18480$
and a prime $d\equiv N_0/2-1\bmod N_0$. Then $d\equiv7\bmod8$, and every prime dividing $N_0$ splits in
$K=\mathbb{Q}(\sqrt{-d})$. \emph{Admissible data} for $P$ are a record $D$ with these components: real scales
$x$, $x_1$, $b$ (stored by their logarithms); natural numbers $\mu$, $m$, $r$: $m$ a power of $2$ with
$\sqrt x\le m\le2\sqrt x$, the exponent $\mu$ with $2^{\mu}=m^{22}$ and $\mu\le32\log x$, and
$r$ a power of $3$ with $x^5\le r\le3x^5$; a finite set $\ovthick(D)$ of infinite places of the field $E(D)$ defined
below; two distinct primes $g_0$, $g_1$ and an integer $\pi_g$ of $K$ with $\pi_g\bar\pi_g=g_0g_1$. The scales satisfy
\[
(\log b)^{100}\le x,\qquad x^{10^6}\le x_1,\qquad x_1^{10^6}\le b,\qquad(\log x_1)^6\le\log b,
\]
and there are further conditions, stated in Chapter~III: a lower bound for $\log\log b$ by a constant of $P$,
congruences and sizes for $g_0$, $g_1$ (both between $b^{10}$ and $b^{11}$), conditions on $\ovthick(D)$, and the
existence of primes in certain short windows.}

So $\log b\ll x\ll x_1\ll b$, with large gaps, and $m\approx\sqrt x$ is the number of theta lifts: the saving
$x^{-1/2}$ of the one-sided bound will be $1/m$. The following objects are functions of $D$.

\emph{The field and its places.} $N_1(D)$ is $N_0\,2^{\mu}=N_0\,m^{22}$ times a power of $3$ fixed by $\mu$, and $E(D)$ is $K$
with the roots of unity of order $N_1(D)$ adjoined: a CM field, abelian over $\mathbb{Q}$. Its infinite places
form a principal homogeneous space under $\operatorname{Gal}(E(D)/K)$. Restriction to the roots of unity identifies
this group with $(\mathbb{Z}/N_1(D))^\times$: it is onto because $K$ is not contained in the field that these roots
of unity generate over $\mathbb{Q}$, the prime $d$ being ramified in $K$ and prime to $N_1(D)$. A base place
identifies the places with this group. So a place $\tau$ has a class modulo $N_1(D)$, and $\nu\cdot\tau$ is defined for
integers $\nu$ prime to $N_1(D)$. The set $\ovtops(D)$ consists of the places of class $\equiv1\bmod N_0$; it has
$N_1(D)/N_0$ elements, and $f(D)=\varphi(N_0)\cdot\#\ovtops(D)$ is the number of all infinite places. The places of
$\ovthick(D)$ have class $\equiv-1\bmod N_0$.

\emph{Sizes.} $k(D)=r+1-m$ and $t(D)=5r$; $a(D)$ is the least multiple of $N_1(D)$ that is at least $x\,b$; and
$\ovbudget(D)=f(D)\,a(D)\,(\log b)^3$.

\emph{Prime sums.} The \emph{nominal weight} of an infinite place $\tau$ is $w_D(\tau)=a(D)/x$ if $\tau\in\ovthick(D)$ and
$w_D(\tau)=a(D)/x_1$ otherwise. For real $z$ and an infinite place $\tau$, with $(\,\cdot\,)_+$ the positive part, put
\begin{gather*}
\ovAsk_D(z,\tau)=\sum_{\nu=1}^{a(D)}\lambda_{d,1}(\nu)\Bigl(\frac z\nu-w_D\bigl((-\nu)\cdot\tau\bigr)\Bigr)_{+},\\
\ovrem(D)=\sum_{\tau\in\ovtops(D)}\Bigl(\ovAsk_D\bigl(a(D),\tau\bigr)-a(D)\bigl(\log a(D)+\kappa_{d,1}-\bar w_D\bigr)\Bigr).
\end{gather*}
Here $(-\nu)\cdot\tau$ is the image of $\tau$ under the class of $-\nu$ modulo $N_1(D)$ when $\nu$ is prime to $N_1(D)$, and
$\tau$ itself otherwise; the terms with $\nu$ not prime to $N_1(D)$ vanish in any case, since every prime dividing $N_1(D)$
splits in $K$ (Setting~1), so that $\Lambda(\nu)=0$ or $1-\chi_K(\nu)=0$. Further, $\bar w_D$ is the mean of $\log w_D(\tau)$ over all infinite
places $\tau$ of $E(D)$. (The formal text writes $\kappa_{d,1}$ as $-1$ minus Euler's constant plus
$\ovRe(L'/L)(1,\chi_K)$; a lemma of Chapter~II identifies the two.) In Section~4 the number $\ovAsk_D(z,\tau)$ will be
the divisibility that the primes not split in $K$ give to a derivative of order $z$ at $\tau$, and $\ovrem(D)$ is
what the tops receive beyond the expected value. Finally the \emph{main amount} is
$\ovhmain(D)=m\,k(D)\,\ovU(D)$, where $\ovU(D)$, the \emph{cost of one direction}, is a sum over all infinite places
$\tau$: the size charged to a derivative of order about $w_D(\tau)$, a factorial against a power of a period, less the
divisibility $\ovAsk_D(w_D(\tau),\tau)$ that such a derivative is sure to have. Chapter~III has the formula. Its value
plays no role here: it cancels.

\ovhead{The remainder is a combination of the remainders of Section 2.} As $\tau$ runs over $\ovtops(D)$, the place
$(-\nu)\cdot\tau$ meets every thick place exactly once if $\nu\equiv1\bmod N_0$, and none otherwise. This is where the
progression $1\bmod N_0$ enters, and it gives the identity
\begin{equation}\label{ov:rem}
\ovrem(D)=a(D)\Bigl(\#\ovthick(D)\,\bigl(R_{d,N_0}(x)-R_{d,N_0}(x_1)\bigr)+\#\ovtops(D)\,R_{d,1}(x_1)\Bigr).
\end{equation}

\ovthm{Main inequality}{In Setting~1, for all parameters $P$ and all admissible data $D$ for $P$, with number of lifts
$m$ and third scale $b$,
\[m\cdot\ovrem(D)\ \le\ (\log b)^2\cdot\ovbudget(D)\ =\ f(D)\,a(D)\,(\log b)^5 .\]}

\ovhead{From the main inequality to the one-sided bound.} Fix $\varepsilon$ and let $x$ be large. Data exist with first
scale $x$, with $\log b=x^{\varepsilon}$, $\log x_1=x^{\varepsilon/6}$, and with
$\#\ovthick(D)\ge\#\ovtops(D)/\log b$ (Chapter~III; the thick set is a construction of Behrend type). Solving
\eqref{ov:rem} for $R_{d,N_0}(x)$ leaves the combination $R_{d,N_0}(x_1)-X\,R_{d,1}(x_1)$ with
$X=\#\ovtops(D)/\#\ovthick(D)\le\log b$. No pointwise prime number theorem in progressions is used to bound it.
Instead, since both remainders have transforms without poles on $\ovRe s\ge1$, every long window contains a good
point: for $v_0=10^6\log x$ large and every $X\in[0,e^{v_0}]$ there is $v_1\in[v_0,v_0^2]$ with
$R_{d,N_0}(e^{v_1})-X\,R_{d,1}(e^{v_1})\le e^{-v_0}$. Replacing $x_1$ by $e^{v_1}$ keeps the data admissible. For the new
data $D'$ the main inequality and \eqref{ov:rem} give
\[
R_{d,N_0}(x)\ \le\ \frac{f(D')\,(\log b)^5}{m\cdot\#\ovthick(D')}+x^{-10^6}
\ \le\ \frac{\varphi(N_0)\,(\log b)^6}{\sqrt x}+x^{-10^6}\ \le\ x^{8\varepsilon-1/2}.
\]

\ovsec{4}{Theorems A, B and C}

One covolume is bounded from both sides. Theorems A and B produce many automorphic forms whose theta lifts are short
and almost integral; Theorem C says that integral forms cannot be that short. The two bounds differ by
$m\cdot\ovrem(D)$.

\ovset{Setting 2 (forms and lifts)}{Let $P$ and $D$ be as in Setting~1. Chapters III, V and VII attach to $D$ the
following objects.

\emph{Levels.} The number $g(D)=g_0g_1$; a prime $N(D)$ that splits completely in $E(D)$, with
$m\log b\le\log N(D)\le30\,m\log b$; an integer $\pi_N(D)$ of $E(D)$ whose product with its complex conjugate is
$N(D)$. The \emph{chosen primes} are
the primes of $E(D)$ that divide $\bar\pi_g$, and $\ovG(D)$ is the product of the groups $\mathrm{GL}_r$ of their
residue fields, a finite group.

\emph{Columns.} $\ovCol(D)$ is a finite-dimensional complex space of algebraic modular forms on the unitary group of a
definite hermitian space of rank $k(D)$ over $E(D)$, of a fixed weight, with level at the chosen primes and at the
primes dividing $\pi_N(D)$. Its subspace $\ovColram(D)$ of \emph{ramified} columns is the image of the product, over
the primes $\mathfrak{p}$ dividing $\pi_N(D)$, of the operators $1-(\text{translation by a central element at }\mathfrak{p})$.
For a prime $\ell$ and a ring homomorphism $j\colon\mathbb{C}\to\mathbb{C}_\ell$, let $O_j$ be the ring of the complex
numbers whose image under $j$ has absolute value at most~$1$. The values of these forms are vectors of a weight
representation, and integrality is measured in a lattice that moves with the point: a form is \emph{$j$-integral} if at
every finite adelic point its value lies in the weight lattice at that point, that is, after $j$ is applied to its
coordinates and then the local weight action of the inverse of the point, all coordinates have absolute value at
most~$1$. A column is
\emph{integral} if it is $j$-integral for every prime $\ell$ and every $j$.

\emph{Rows.} $\ovRow(D)$ is a finite-dimensional complex space of cusp forms on the unitary group of an indefinite
hermitian space of rank $r$ over $E(D)$: of a fixed discrete series type at infinity, of principal level at the chosen
primes and with no level at $N(D)$, cut out by a fixed idempotent of the group algebra of $\ovG(D)$.

\emph{Poincaré series.} For $\gamma\in\ovG(D)$, the linear map $\ovP_D(\gamma)\colon\ovColram(D)\to\ovRow(D)$ multiplies a
ramified column by a fixed non-zero constant, forms its Poincaré series on the indefinite group, translates it by
$\gamma$ and applies $\#\ovG(D)$ times the idempotent.

\emph{Twists and lifts.} An \emph{$m$-fold twist} of $D$ is a continuous character $\vartheta$ of the norm-one adeles of
$E(D)$, trivial on the norm-one elements of $E(D)$ and at the infinite places, at most tamely ramified, such that the
set $Q(\vartheta)$ of primes where it ramifies contains the complex conjugate of none of its members, consists of
primes of norm prime to $g(D)N(D)$, and is also the set of ramified primes of each of $\vartheta^2,\dots,\vartheta^m$. The trivial
character $1$ is one. For an $m$-fold twist $\vartheta$, the map
$\Theta_{D,\vartheta}\colon\ovRow(D)\to\ovA_{D,\vartheta}$ is the composite of $m$ theta lifts, the $i$-th formed with
$\vartheta^{i}$; it is injective, and $\ovA_{D,\vartheta}$ is a space of algebraic modular forms on a definite unitary
group of rank $r+m\,t(D)$. For $F,F'\in\ovA_{D,\vartheta}$ and $\sigma\in\ovAut$, $\langle F,F'\rangle_\sigma$ is the
Petersson product, a finite sum over classes, of the forms whose values are the $\sigma$-conjugates of those of $F$ and
$F'$. Put $\|F\|_\sigma^2=\langle F,F\rangle_\sigma$ and
$\ovcov_\sigma^2(F_1,\dots,F_n)=|\det(\langle F_i,F_{i'}\rangle_\sigma)_{i,i'}|$. For forms in $\ovA_{D,\vartheta}$, $j$-integrality is defined as for columns, in the weight lattice at each point.

\emph{Index.} Rows $y_1,\dots,y_n$ have \emph{$D$-index at most $\delta$} if there are real numbers $\eta_\ell$,
$\ell\in\mathbb{N}$, with $\sum_{\ell\in J}\eta_\ell\le\delta$ for every finite $J\subset\mathbb{N}$, such that for every
prime $\ell$, every $j\colon\mathbb{C}\to\mathbb{C}_\ell$ and every $m$-fold twist $\vartheta$ of $D$ with $\ell$ prime to
the norms of the primes in $Q(\vartheta)$ there is an $n\times n$ matrix $u$ over $O_j$ with
$|j(\det u)|_\ell\ge e^{-\eta_\ell}$ for which $\Theta_{D,\vartheta}\bigl(\sum_{i'}u_{ii'}\,y_{i'}\bigr)$ is $j$-integral
for $i=1,\dots,n$.}

\ovthm{Theorem A (short forms)}{In Settings 1 and 2 there is a rational number $q>0$ with the following two properties.

\emph{(a) Size.} There is a set $S$ of integral ramified columns that spans $\ovColram(D)$ over $\mathbb{C}$ such that,
for all $\psi\in S$, all $\gamma\in\ovG(D)$ and all $\sigma\in\ovAut$,
\[\bigl\|\Theta_{D,1}\bigl(q\,\ovP_D(\gamma)\psi\bigr)\bigr\|_\sigma^{2}\ \le\ \exp\bigl(\ovhmain(D)-m\cdot\ovrem(D)+\log b\cdot\ovbudget(D)\bigr).\]

\emph{(b) Index.} For every finite set $\ovT\subseteq\ovG(D)$, every $n\ge0$ and all rows $y_1,\dots,y_n$, each of the form
$q\,\ovP_D(\gamma)\psi$ with $\gamma\in\ovT$ and $\psi$ an integral ramified column, the family $y_1,\dots,y_n$ has $D$-index
at most $\#\ovT\cdot\dim\ovCol(D)\cdot\ovbudget(D)$.}

\emph{Idea (the authors' comments; Chapter V).} A seesaw identity moves the problem to definite groups: for every
twist, the chain of lifts of the Poincaré series of a ramified column is a chain of $m$ theta lifts
of the column itself between definite spaces, restricted at the end. There a lift is a finite sum of Taylor
coefficients at the origin of theta functions with complex multiplication by $K$, of order $a(D)$ at the tops. Their
size at every conjugate comes from Stirling's formula and a period. Their divisibility at a prime $p$ not split in $K$
comes from the formal group of the CM curve, of height two at $p$: a derivative of order $z$ against a partner place of
weight $w$ gains about $(z/p-w)_+$ powers of $p$. Summed over $p$ with weights $\log p$ this is $\ovAsk_D$, which is why
$\lambda_{d,M}$ appears. The rational $q$ divides the gain out, a power of the prime $N(D)$, chosen in a short window,
pays the excess of size over expected divisibility, and $\ovhmain(D)-m\cdot\ovrem(D)$ is left. What
some primes lose is regained on the image of scaled Hecke operators, at the price of the $D$-index in (b).

\ovthm{Theorem B (enough forms)}{In Settings 1 and 2, $\dim\ovCol(D)>0$ and $\dim\ovRow(D)>0$, and there is a finite set
$\ovT\subseteq\ovG(D)$ with
\begin{gather*}
\#\ovT\cdot\dim\ovCol(D)\ \le\ 5\,\dim\ovRow(D)\cdot\log(1+b^2)\qquad\text{and}\\
\dim\sum_{\gamma\in\ovT}\ovP_D(\gamma)\bigl(\ovColram(D)\bigr)\ \ge\ \Bigl(1-\frac4{b^2}\Bigr)\dim\ovRow(D).
\end{gather*}}

\emph{Idea (Chapter VI).} The translates are spread: for every subspace $W$ of $\ovRow(D)$ some $\gamma$ pulls $W$ back
to at most the share $\tfrac12\bigl(1+\dim W/\dim\ovRow(D)\bigr)+b^{-2}$ of the columns. This follows from a second
moment of the Poincaré series over $\ovG(D)$, of Petersson type, whose main term consists of Bessel integrals at the
real places (the place $A=\mathbb{C}$ of Chapter VI, which the notes of that chapter call the complex place) and of
orbital integrals at the chosen primes, and from a variance inequality. A greedy choice of
translates then gains at each step in proportion to the dimension still missing, whence the logarithm.

\ovthm{Theorem C (integral forms are long)}{In Settings 1 and 2, let $n\ge0$ be an integer, let
$y_1,\dots,y_n$ be rows in $\ovRow(D)$ that are linearly independent over $\mathbb{C}$, and let $\delta$ and $V$ be real
numbers. Suppose that $n\ge(1-4/b^2)\dim\ovRow(D)$, that $y_1,\dots,y_n$ have $D$-index at most $\delta$, and that
$\ovcov_\sigma^2\bigl(\Theta_{D,1}y_1,\dots,\Theta_{D,1}y_n\bigr)\le e^{V}$ for all $\sigma$ in $\ovAut$. Then
\[\dim\ovRow(D)\cdot\bigl(\ovhmain(D)-3\log b\cdot\ovbudget(D)\bigr)\ \le\ V+2\delta .\]}

\emph{Idea (Chapters IV, VII and VIII).} The tool is a height for families $Z=(Z_{\ell,\iota})$ of complex numbers
indexed by the pairs of a prime $\ell$ and a field isomorphism $\iota\colon\mathbb{C}\to\bar{\mathbb{Q}}_\ell$. Such a
family has \emph{height at most $h$} if there are $H\in\mathbb{C}^\times$, $\alpha\in\mathbb{R}$ and real numbers
$\eta_\ell$ with $|\sigma(H)|\le e^{\alpha}$ for all $\sigma\in\ovAut$, $|\iota(H\,Z_{\ell,\iota})|_\ell\le e^{\eta_\ell}$ for
all $(\ell,\iota)$, and such that $\alpha$ plus the sum of the $\eta_\ell$ over the primes $\ell$ below any sufficiently
large bound is at most $h$. Heights add in
products, and a height bound for the constant family $1$ is non-negative: this is the product formula, applied to the
norm of $H$. Let $\mathfrak{z}_D(\vartheta)=\det(\Theta_{D,\vartheta}^{*}\Theta_{D,\vartheta})$. For a fit family of
twists $\theta$ (Section~5), two bounds hold. First, $\mathfrak{z}_D(\theta_{\ell,\iota})/\mathfrak{z}_D(1)$ has height at
most $V+2\delta+b\,\dim\ovRow(D)$: the family is completed to a basis by crude rows, $V$ bounds the Gram determinant at
the trivial twist at every conjugate, and the $D$-index bounds the denominators at each $\ell$. Second,
$\mathfrak{z}_D(1)^2/\mathfrak{z}_D(\theta_{\ell,\iota})^2$ has height at most
$2\dim\ovRow(D)\,(2\log b\cdot\ovbudget(D)-\ovhmain(D))$: the volume of the chain is a constant times the volumes of the
$m$ single lifts, a single lift is governed by Hecke $L$-values, and each ratio of an untwisted to a twisted value
repays one direction, $\ovU(D)$. The square of the first family times the second is the constant family $1$, so
twice the first bound plus the second is non-negative. This is Theorem C, since $b\le\log b\cdot\ovbudget(D)$.

\ovhead{A, B and C give the main inequality.} Take $q$ and $S$ from A and $\ovT$ from B. The rows
$q\,\ovP_D(\gamma)\psi$, $\gamma\in\ovT$, $\psi\in S$, span $\sum_{\gamma\in\ovT}\ovP_D(\gamma)(\ovColram(D))$; choose a
basis $y_1,\dots,y_n$ of this space among them, so that $n\ge(1-4/b^2)\dim\ovRow(D)$. By Hadamard's inequality and
A(a), $\ovcov_\sigma^2(\Theta_{D,1}y_1,\dots,\Theta_{D,1}y_n)\le e^{V}$ with
$V=n\,(\ovhmain(D)-m\,\ovrem(D)+\log b\cdot\ovbudget(D))$. By A(b) and B their $D$-index is at most
$\delta=5\dim\ovRow(D)\log(1+b^2)\,\ovbudget(D)$. Theorem C compares the two. A crude bound,
$|\ovhmain(D)|+m\,|\ovrem(D)|\le b^2x^{99}$, shows that replacing $n$ by $\dim\ovRow(D)$ costs at most $4x^{99}$ for each
row. Dividing by $\dim\ovRow(D)$, the main amount cancels:
\[
m\cdot\ovrem(D)\ \le\ \bigl(4\log b+10\log(1+b^2)\bigr)\,\ovbudget(D)+4x^{99}\ \le\ (\log b)^2\,\ovbudget(D).
\]

\ovsec{5}{The twists}

Theorem C compares the lifts at the trivial twist with the lifts at twists where they have few congruences with
other forms. The formal proof does not find one such twist for all primes: it finds one for each embedding, which
is what the height of Section~4 allows.

\ovstm{Twists}{In Settings 1 and 2, for all parameters $P$ and all admissible data $D$ for $P$ there is a family
$\theta=(\theta_{\ell,\iota})$ of $m$-fold twists of $D$, one for each prime $\ell$ and each field isomorphism
$\iota\colon\mathbb{C}\to\bar{\mathbb{Q}}_\ell$, that is \emph{fit for $D$}.}

Fitness is defined in Chapter~VIII; it has three clauses. The prime $\ell$ is prime to the norms of the primes in
$Q(\theta_{\ell,\iota})$. For each $i$ with $1\le i\le m$ (Chapter~VIII, like the formal text, calls this index $j$),
the determinants that measure the $\ell$-adic denominators of the orthogonal projection of integral forms to the
space $Z_i$ (the image, in the forms of the auxiliary pair, of the adjoint of the $i$-th lift at
$\theta_{\ell,\iota}^{\,i}$ restricted to good forms), and to certain lines of lifts from $\mathrm{U}(1)$, are at most
$e^{r'\eta_\ell}$ for each dimension, where $r'$ is the rank of the auxiliary space ($r(D)$ for the good lifts, $1$ for
the lines), with real numbers $\eta_\ell$
whose sum over any finite set of primes is at most $\ovbudget(D)/(m\,k(D))$. And for finitely many Hecke characters
of $\mathrm{U}(1)$, the inverse ratios of twisted to untwisted $L$-values have height at most $x^{999}$.

\emph{Idea (the authors' comments; Chapter VIII).} A denominator of the projection to theta lifts is a congruence
between a lift and a form that is not a lift. By Ribet's method, in a form with congruent blocks and in any rank, such
a congruence produces classes in Selmer groups of the Galois representations attached to the forms (attached in the
weak sense of Chapter~VII); these
representations are an irreducible block plus a string of powers of the cyclotomic character. The even classes of
these groups, where in the authors' words Vandiver's conjecture lives, are never bounded. Twisting by an
anticyclotomic character of $\ell$-power order, ramified at new primes where a tie between inertia and Frobenius is
imposed (an intermediate device, which the final twist releases), makes all the relevant Selmer groups short for the first $m$ powers at once; Taylor--Wiles primes in a
triangular arrangement remove the even classes, the odd parts of Tate--Shafarevich groups of abelian CM fields
are bounded by one explicit integer (for the untwisted character the module is the unramified part, with no oddness
condition), and the conditions at $\ell$ cost a loss that is controlled by lifting crystalline extension classes in
large weight.

\ovsec{6}{What is quoted}

The repository separates the argument from a background library, which proves from Mathlib the known results
that the argument needs, often in a weaker form that suffices. This write-up reads the argument and quotes the
background. That is a choice of exposition, not an extra hypothesis: the formal proof contains proofs of all of it.
The quoted results come from \emph{analytic number theory} (Landau's theorem, the zero-free box for quadratic
twists that is discussed below, a weak form of the Siegel--Walfisz theorem; from Mathlib, the functional equation and the non-vanishing on
$\ovRe s=1$); from \emph{class field theory and Galois cohomology} (Artin reciprocity, Chebotarev's theorem, local
Tate duality, the Greenberg--Wiles formula); from \emph{complex multiplication} (theta functions of CM curves and
their formal groups, Kronecker's limit formula, the Chowla--Selberg formula); from \emph{automorphic forms}
(reduction theory, the Weil representation and theta lifts, the Siegel--Weil and Rallis inner product formulas,
discrete series); and from \emph{Galois representations} (those of eigenforms on unitary groups, through Shimura
varieties; $p$-adic Hodge theory with the Bloch--Kato conditions and Kisin modules; bounds of Kolyvagin--Rubin
type). The catalogue at the end of this handout lists each result with a precise statement, its reference where the
place is confirmed, and the relation of the formal statement to the published one: the theorem itself, a special
case, or a weaker form.

Some quoted results stand outside these three relations: no published source in the form needed has been
identified, so they are quoted from the formal library, with their statements in the chapters, and not from the
literature. The first is the zero-free box of Section~2, one prime $d$ in a class
whose twists $\tilde L(s,\xi\chi_K)$ have no zero in a box for all characters $\xi$ modulo $M$ at once. The formal
library proves it in two steps. The first is an estimate of zero-density type for each fixed character $\xi$,
real or not: the number of the primes $d\le x$ with $d\equiv7\bmod8$, prime to $M$, for which
$\tilde L(\,\cdot\,,\xi\chi_K)$, $K=\mathbb{Q}(\sqrt{-d})$, has a zero in a given box to the right of the critical
line is at most a constant times a power of $x$ with exponent less than $1$. The second adds these bounds over the characters $\xi$ and sets
them against the divergence of the sum of $1/p$ over the primes $p$ of a residue class. The source that the formal
text names, a zero-density theorem of Jutila, concerns real characters. It does not give the first step for a
non-real $\xi$, whose twists are not real, and no published place for that step has been confirmed; its statement is
in Chapter~II. The others are in Chapter~VIII: a bound for the odd parts of the Tate--Shafarevich groups of an
abelian CM field by one explicit integer, uniform in the prime, the twist and the precision (Greither's theorem is an
identity of characteristic ideals, and how it would give such a bound is not established); a finite set of Galois
elements that detects the relevant cocycles; and the relations at split primes, with a lower bound for the number of
relation families, together with the count of trivialised cocycles that supplies its hypothesis.

\Needspace{19\baselineskip}
\ovsec{7}{A map}

{\small\setlength{\tabcolsep}{4pt}\noindent
\begin{tabular}{@{}>{\raggedright\arraybackslash}p{8mm}>{\raggedright\arraybackslash}p{90mm}>{\raggedright\arraybackslash}p{46mm}@{}}
\toprule\textbf{ch.} & \textbf{what it proves} & \textbf{units (folders)}\\\midrule
I & The five theorems from (3); the main inequality from A, B, C; the one-sided bound from the main inequality &
  statement file, main file, head file of the argument\\
II & The remainder and its residues; the one-sided criterion; a good point in every window; identity
  \eqref{ov:rem}; the last steps of A, B and C & Endgame\\
III & Parameters and data; the field $E(D)$ and its places; comparison of sizes; existence of data at every large
  $x$ & Data\\
IV & Heights and the product formula; few translates nearly span; the height of one Gram determinant & Heights\\
V & Theorem A: the seesaw identity, divisibility of derivatives of theta functions, sizes, the level & Seesaw, Curve,
  Columns, TwistedHecke, Core\\
VI & Theorem B: the second moment of Poincaré series, Bessel functions at real places & Filling, RealPlace,
  DiscreteSeries\\
VII & Theorem C: volume of the chain as a product of lift volumes, the exchange identity, heights of ratios of
  $L$-values & Rows, Exchange, Trace, Forms, SiegelWeil, Weil, GaloisRep; DiscreteSeries (Theorems 2, 3, 6)\\
VIII & The twists: Ribet's method, Selmer groups under twisting, Taylor--Wiles primes, crystalline lifts & Selmer,
  Patching, Lift\\
\bottomrule
\end{tabular}\par}

\Needspace{9\baselineskip}
\ovhead{Where each statement lives.} Paths are relative to the source folder of the compact edition at commit
\code{\PinShort}; the numbers are the lines of the declaration. The hypotheses of Theorem C stand in the
\code{variable} line 13--14 of its file.

{\footnotesize\setlength{\tabcolsep}{3pt}
\begin{longtable}{@{}>{\raggedright\arraybackslash}p{43mm}>{\raggedright\arraybackslash}p{57mm}>{\raggedright\arraybackslash}p{45mm}@{}}
\toprule\textbf{statement} & \textbf{Lean name} & \textbf{file, lines}\\\midrule\endhead
%%ROWS-BEGIN (made by a program from reading/census.json; do not type)
Theorems (1), (2), (3), (4), (5) & \code{riemannHypothesis}\newline \code{riemannZeta\_\allowbreak eq\_\allowbreak zero\_\allowbreak re}\newline \code{LFunction\_\allowbreak ne\_\allowbreak zero\_\allowbreak of\_\allowbreak half\_\allowbreak lt\_\allowbreak re}\newline \code{LFunction\_\allowbreak eq\_\allowbreak zero\_\allowbreak re}\newline \code{LFunction\_\allowbreak eq\_\allowbreak zero\_\allowbreak re\_\allowbreak of\_\allowbreak isPrimitive} & \code{Main.lean}, 4, 6, 8, 10, 12--13\\[3pt]
The propositions (2) to (5) & the definitions of the statement file & \code{Statements.lean}, 3--16\\[3pt]
(3) inside the argument; from (3) to the line, for characters and for $\zeta$ & \code{New.\allowbreak grh}\newline \code{New.\allowbreak zeros\_\allowbreak on\_\allowbreak line}\newline \code{New.\allowbreak zeta\_\allowbreak zeros\_\allowbreak on\_\allowbreak line} & \code{New/\allowbreak Main.lean}, 174--190, 194--196, 200--204\\[3pt]
The weight $\lambda_{d,M}$; the remainder $R_{d,M}$; zeros as poles & \code{lambdaClass}\newline \code{primeError}\newline \code{isPolar\_\allowbreak primeError} & \code{New/\allowbreak Endgame/\allowbreak PrimeError.lean}, 25, 38--39, 69--111\\[3pt]
One-sided criterion; a good point in every window & \code{entireL\_\allowbreak ne\_\allowbreak zero\_\allowbreak of\_\allowbreak primeError\_\allowbreak le}\newline \code{IsPolar.\allowbreak exists\_\allowbreak window\_\allowbreak sub\_\allowbreak mul\_\allowbreak le} & \code{New/\allowbreak Endgame/\allowbreak OneSidedCriterion.lean}, 28--36, 73--99\\[3pt]
Landau's theorem (quoted) & \code{IsPolar.\allowbreak landau} & \code{Lib/\allowbreak Analytic/\allowbreak Landau/\allowbreak PolarLandau.lean}, 88--136\\[3pt]
Zero-free twists for one prime (quoted) & \code{exists\_\allowbreak entireL\_\allowbreak quadTwist\_\allowbreak ne\_\allowbreak zero} & \code{Lib/\allowbreak Analytic/\allowbreak QuadraticTwist/\allowbreak ZeroFreeTwist.lean}, 193--206\\[3pt]
One-sided bound; main inequality & \code{eventually\_\allowbreak primeError\_\allowbreak le\_\allowbreak rpow}\newline \code{star\_\allowbreak inequality} & \code{New/\allowbreak Main.lean}, 149--167, 121--145\\[3pt]
Setting 1: parameters; data; their existence & \code{Parameters}\newline \code{Data}\newline \code{exists\_\allowbreak data} & \code{New/\allowbreak Data/\allowbreak Different.lean}, 23--28\newline \code{New/\allowbreak Data/\allowbreak AdmissibleData.lean}, 22--45\newline \code{New/\allowbreak Data/\allowbreak DataExist.lean}, 49--108\\[3pt]
The main amount; the remainder; the crude bound; identity (1) & \code{Data.\allowbreak mainHeight}\newline \code{Data.\allowbreak remainder}\newline \code{Data.\allowbreak abs\_\allowbreak mainHeight\_\allowbreak add\_\allowbreak remainder\_\allowbreak le}\newline \code{Data.\allowbreak remainder\_\allowbreak eq} & \code{New/\allowbreak Endgame/\allowbreak ShortAndEnough.lean}, 32, 42, 74--82, 86--112\\[3pt]
Setting 2: rows; columns; the chain of lifts & \code{Data.\allowbreak rows}\newline \code{Data.\allowbreak Column}\newline \code{Data.\allowbreak Θ} & \code{New/\allowbreak Trace/\allowbreak GoodIdempotent.lean}, 118\newline \code{New/\allowbreak Columns/\allowbreak ColumnDefs.lean}, 71\newline \code{New/\allowbreak Rows/\allowbreak RowVolume.lean}, 27\\[3pt]
Poincar\'e maps; integral ramified columns; index & \code{Data.\allowbreak projectedPoincare}\newline \code{Data.\allowbreak integralRamified}\newline \code{Data.\allowbreak IndexLe} & \code{New/\allowbreak Endgame/\allowbreak ShortAndEnough.lean}, 120--121, 125--126, 130--132\\[3pt]
Theorems A and B & \code{Data.\allowbreak short\_\allowbreak forms}\newline \code{Data.\allowbreak enough\_\allowbreak forms} & \code{New/\allowbreak Endgame/\allowbreak ShortAndEnough.lean}, 138--168, 173--178\\[3pt]
Theorem C; the family given by A and B & \code{Data.\allowbreak integral\_\allowbreak forms\_\allowbreak long}\newline \code{Data.\allowbreak exists\_\allowbreak family\_\allowbreak covolumeSq\_\allowbreak le} & \code{New/\allowbreak Endgame/\allowbreak LongForms.lean}, 63--71, 78--92\\[3pt]
Height; the product formula & \code{HeightLe}\newline \code{HeightLe.\allowbreak nonneg} & \code{New/\allowbreak Heights/\allowbreak Heights.lean}, 21--23, 65--77\\[3pt]
Fit families of twists; their existence & \code{Data.\allowbreak IsFitFamily}\newline \code{exists\_\allowbreak isFitFamily} & \code{New/\allowbreak Selmer/\allowbreak FitTwists.lean}, 256--263\newline \code{New/\allowbreak Selmer/\allowbreak FitTwistsExist.lean}, 24--65\\[3pt]
%%ROWS-END
\bottomrule
\end{longtable}}
